% Generated draft: requires mathematical and source review.
\documentclass[11pt]{article}
\usepackage[margin=25mm]{geometry}
\usepackage{fontspec}
\setmainfont{Latin Modern Roman}
\usepackage{amsmath,amssymb,mathtools,unicode-math}
\setmathfont{Latin Modern Math}
\usepackage{enumitem,needspace,xurl,hyperref}
\setlength{\parindent}{0pt}
\setlength{\parskip}{5pt}
\setlength{\emergencystretch}{3em}
\newcommand{\N}{\mathbb{N}}
\newcommand{\Z}{\mathbb{Z}}
\newcommand{\Q}{\mathbb{Q}}
\newcommand{\R}{\mathbb{R}}
\newcommand{\C}{\mathbb{C}}
\newcommand{\F}{\mathbb{F}}
\newcommand{\A}{\mathbb{A}}
\newcommand{\PP}{\mathbb{P}}
\newcommand{\E}{\mathbb{E}}
\newcommand{\eps}{\varepsilon}
\newcommand{\dd}{\,\mathrm{d}}
\begin{document}
\section*{UnsolvedMath ID: 701418}
\textbf{Source problem code:} catalogue entry
\subsection{Definitions and mathematical statement}
Work in the first-order language of groups \(\mathcal L_{\mathrm{grp}}=\{\cdot,{}^{-1},1\}\). For a group \(G\), its elementary theory is
\[
\operatorname{Th}(G)=\{\varphi:\varphi\text{ is an }\mathcal L_{\mathrm{grp}}\text{-sentence and }G\models\varphi\}.
\]
This theory is decidable if there exists an algorithm that, on input any \(\mathcal L_{\mathrm{grp}}\)-sentence \(\varphi\), halts and correctly determines whether \(G\models\varphi\).

For groups \(A\) and \(B\), their free product \(A*B\) is the group equipped with canonical homomorphisms \(i_A:A\to A*B\) and \(i_B:B\to A*B\) such that, for every group \(H\) and homomorphisms \(f:A\to H\) and \(g:B\to H\), there exists a unique homomorphism \(h:A*B\to H\) satisfying \(h\circ i_A=f\) and \(h\circ i_B=g\).

\textbf{Question (O. Kharlampovich; described in the source as a well-known problem).} For every pair of groups \(A,B\), does
\[
\bigl(\operatorname{Th}(A)\text{ is decidable}\bigr)\ \land\ \bigl(\operatorname{Th}(B)\text{ is decidable}\bigr)
\quad\Longrightarrow\quad
\operatorname{Th}(A*B)\text{ is decidable}
\]
hold? No finiteness, generation, presentation, or nontriviality assumptions on the factors are stated in the supplied record.
\subsection{Short English statement}
If two arbitrary groups each have a decidable first-order theory in the language of groups, must their free product also have a decidable first-order theory?
\subsection{Sources}
\begin{itemize}
\item Local UnsolvedMath registry, permanent ID \(701418\), exact source problem code \texttt{catalogue entry}; supplied record titled \emph{Kourovka Notebook Problem 17.117}, upstream revision \texttt{d371ee917ac9647fbb06cd6fc529deec6f3dda82}.
\item \emph{The Kourovka Notebook (Unsolved Problems in Group Theory)}, 17th issue (2010), Problem 17.117. The supplied record attributes the proposal to O. Kharlampovich and calls it a well-known problem.
\item The supplied background also identifies the 21st issue edition (2026), arXiv:1401.0300v48. The collection description names E. I. Khukhro and V. D. Mazurov as editors. These bibliographic details are supplied attribution and have not been independently verified.
\item The local record labels the problem \texttt{open} and gives an update date of 2026-10-16. This is supplied registry metadata, not an independently verified assessment of the literature.
\end{itemize}
\end{document}
